%% notation.tex --- the notation table, front matter
%% Symbols here match the companion lecture slides. Keep the two in step.

\chapter*{Notation}
\addcontentsline{toc}{chapter}{Notation}
\markboth{Notation}{Notation}

One set of symbols is used throughout the book. Where a chapter needs more, it
extends this table rather than redefining it.

\vspace{4pt}
\noindent\textbf{Data.}
\begin{center}
\small
\renewcommand{\arraystretch}{1.25}
\begin{tabular}{@{}L{2.6cm}L{7.4cm}L{1.6cm}@{}}
\toprule
\textbf{Symbol} & \textbf{Meaning} & \textbf{From}\\
\midrule
$x$                 & an input variable, or \emph{feature} & Ch.\ 1\\
$y$                 & the output, or \emph{target}, we wish to predict & Ch.\ 1\\
$m$                 & number of training examples & Ch.\ 1\\
$n$                 & number of features & Ch.\ 4\\
$\xii{i},\ \yii{i}$ & the $i$-th training example; the superscript is an index, not a power & Ch.\ 1\\
$\xjj{i}{j}$        & feature $j$ of training example $i$ & Ch.\ 4\\
$\vx$               & a feature vector, $\vx\in\Real{n}$ & Ch.\ 4\\
$\Xmat$             & the design matrix, $m\times n$ & Ch.\ 4\\
$\vy$               & the vector of targets, $\vy\in\Real{m}$ & Ch.\ 4\\
$\Dset$             & the training set $\{(\xii{i},\yii{i})\}_{i=1}^{m}$ & Ch.\ 1\\
$\Xspace,\ \Yspace$ & the input space and the output space & Ch.\ 1\\
\bottomrule
\end{tabular}
\end{center}

\vspace{6pt}
\noindent\textbf{Models and learning.}
\begin{center}
\small
\renewcommand{\arraystretch}{1.25}
\begin{tabular}{@{}L{2.6cm}L{7.4cm}L{1.6cm}@{}}
\toprule
\textbf{Symbol} & \textbf{Meaning} & \textbf{From}\\
\midrule
$f$                 & the model, or hypothesis, $f:\Xspace\to\Yspace$ & Ch.\ 1\\
$\Hspace$           & the hypothesis space: every model the algorithm may return & Ch.\ 1\\
$\yhat$             & a prediction, $\yhat=f(x)$ & Ch.\ 1\\
$w,\ b$             & the weight and the bias of a univariate linear model & Ch.\ 2\\
$\vw$               & the weight vector, $\vw\in\Real{n}$ & Ch.\ 4\\
$\fwb,\ \fvwb$      & the model with parameters $(w,b)$, or $(\vw,b)$ & Ch.\ 2\\
$J(w,b)$            & the cost: the objective minimised over the training set & Ch.\ 2\\
$\loss$             & the loss on a single example & Ch.\ 6\\
$\alpha$            & the learning rate, or step size & Ch.\ 3\\
$\lambda$           & the regularization strength & Ch.\ 7\\
$\grad J$           & the gradient of $J$ & Ch.\ 4\\
$\argmin$           & the argument at which a function attains its minimum & Ch.\ 2\\
$\gsig(z)$          & the sigmoid, $\gsig(z)=1/(1+e^{-z})$ & Ch.\ 6\\
$\aL{l}{j}$         & activation of unit $j$ in layer $l$ & Ch.\ 8\\
$\Wl{l},\ \bl{l}$   & weights and biases of layer $l$ & Ch.\ 8\\
\bottomrule
\end{tabular}
\end{center}

\vspace{6pt}
\noindent\textbf{Mathematics.}
\begin{center}
\small
\renewcommand{\arraystretch}{1.25}
\begin{tabular}{@{}L{2.6cm}L{7.4cm}L{1.6cm}@{}}
\toprule
\textbf{Symbol} & \textbf{Meaning} & \textbf{From}\\
\midrule
$\R,\ \Real{n}$     & the reals; the space of real $n$-vectors & Ch.\ 1\\
$\Prob(A)$          & the probability of an event & Ch.\ 6\\
$\Exp[X]$           & the expectation of a random variable & Ch.\ 10\\
$\|\vx\|$           & the Euclidean norm of $\vx$ & App.\ A\\
$\vw\cdot\vx$       & the dot product $\sum_j w_j x_j$ & Ch.\ 4\\
$\Xmat^{\mathsf{T}}$ & the transpose of a matrix & Ch.\ 4\\
$:=$                & assignment, as opposed to the assertion $=$ & Ch.\ 3\\
\bottomrule
\end{tabular}
\end{center}

\vspace{6pt}
\noindent A superscript in parentheses is always an example index, never an
exponent: $\xii{2}$ is the second training example, while $x^2$ is $x$ squared.
Vectors are column vectors and are set in bold, $\vx$; scalars are italic, $x$.

\cleardoublepage
